\documentclass[10pt,twocolumn]{article}
\usepackage[margin=0.65in]{geometry}
\usepackage{amsmath,amssymb,booktabs,array,xcolor}
\usepackage[hidelinks]{hyperref}

\newcommand{\ket}[1]{\left|#1\right\rangle}
\newcommand{\bra}[1]{\left\langle#1\right|}
\newcommand{\braket}[2]{\left\langle#1\middle|#2\right\rangle}
\newcommand{\matel}[3]{\left\langle#1\middle|#2\middle|#3\right\rangle}
\newcommand{\HH}{\mathrm{HH}}
\newcommand{\QSE}{\mathrm{QSE}}
\newcommand{\ED}{\mathrm{ED}}
\newcommand{\Full}{\mathrm{Fullmeta}}
\newcommand{\Geom}{\mathrm{Geom}}
\newcommand{\Span}{\operatorname{span}}
\newcommand{\cond}{\operatorname{cond}}

\definecolor{boxgray}{RGB}{245,245,245}
\newenvironment{paperbox}{\par\smallskip\noindent\begin{center}\begin{minipage}{0.94\linewidth}\small\hrule\smallskip}{\smallskip\hrule\end{minipage}\end{center}\smallskip}

\title{Pedagogical duplicate for Paper III: geometry-selected QSE spectra}
\author{Internal learning source; not a journal draft}
\date{May 31, 2026}

\begin{document}
\maketitle

\section*{Purpose}
This note is the learning version of the Paper-III static path.  The journal draft can stay compact.  This duplicate explains what the QSE calculation is doing, why operator sources matter, why overlap conditioning matters, and why a spectrum can look accurate while the sorted excited-state ladder looks poor.

The reader is assumed to understand Paper I: a ground-state adaptive ansatz is built from an operator sequence and optimized for a Hubbard--Holstein Hamiltonian.  Paper III starts from that prepared ground-state vector and asks which excitation directions should be used to build a spectral subspace.

\section{Paper-I object reused by Paper III}
Paper I gives an optimized ground-state ansatz
\begin{equation}
  \ket{\psi_0}
  = U_m(\theta_m)\cdots U_1(\theta_1)\ket{\phi_{\rm ref}},
  \qquad
  U_j(\theta_j)=e^{-i\theta_j A_j}.
\end{equation}
The sequence
\begin{equation}
  \mathcal A=(A_1,A_2,\ldots,A_m)
\end{equation}
is useful because it records which directions were needed to lower the ground-state energy.  Paper III uses \(\ket{\psi_0}\) as a prepared state.  Some Paper-III baselines also reuse the operator information in \(\mathcal A\) as a candidate excitation source.

\begin{paperbox}
\textbf{Translation from Paper I.}  Paper I studies how to acquire the ground-state ansatz.  Paper III takes the acquired state as fixed input and studies how to acquire a nonorthogonal excitation basis around that state.
\end{paperbox}

\section{Regular QSE as a subspace solve}
Choose excitation operators
\begin{equation}
  \Omega_1,\Omega_2,\ldots,\Omega_M.
\end{equation}
Each operator creates a candidate excited-state direction from the prepared state:
\begin{equation}
  \ket{\phi_a}=\Omega_a\ket{\psi_0}.
\end{equation}
For the Paper-III static mode we usually remove the component parallel to the prepared state,
\begin{equation}
  Q_0 = I-\ket{\psi_0}\bra{\psi_0},
  \qquad
  \ket{\tilde\phi_a}=Q_0\Omega_a\ket{\psi_0}.
\end{equation}
This makes the QSE basis focus on directions away from the reference state.  The nonorthogonal overlap and Hamiltonian matrices are
\begin{align}
  S_{ab} &= \braket{\tilde\phi_a}{\tilde\phi_b},\\
  H_{ab} &= \matel{\tilde\phi_a}{H}{\tilde\phi_b}.
\end{align}
The QSE roots are obtained from the generalized eigenproblem
\begin{equation}
  Hc_\nu = E_\nu Sc_\nu.
  \label{eq:qse_gep}
\end{equation}
The coefficient vector \(c_\nu\) synthesizes a root inside the nonorthogonal basis:
\begin{equation}
  \ket{\Psi_\nu^{\QSE}}
  =\sum_a c_{\nu a}\ket{\tilde\phi_a}.
\end{equation}

\begin{paperbox}
\textbf{Key point.}  QSE is the matrix construction and generalized eigenproblem in Eq.~\eqref{eq:qse_gep}.  A concrete QSE calculation is not fully specified until the excitation operators \(\Omega_a\) are specified.
\end{paperbox}

\section{Operator source versus selection rule}
Two choices are being studied at once.

\begin{enumerate}
\item \textbf{Operator source:} where the candidate operators come from.
\item \textbf{Selection rule:} which candidates are retained in the final QSE basis.
\end{enumerate}
For the weak--weak Hubbard--Holstein pilot, the current rows mean:
\begin{center}
\begin{tabular}{p{0.42\linewidth}p{0.48\linewidth}}
\toprule
Row & Meaning\\
\midrule
regular mixed Fullmeta QSE & Use the full mixed fermion--boson Fullmeta alphabet as the QSE basis.\\
geometry-selected mixed Fullmeta QSE & Start from the same Fullmeta alphabet, then select a compact basis by geometry/conditioning/resource scores.\\
Hamiltonian-term QSE & Use operators derived from Hamiltonian Pauli terms.\\
ground-ansatz-derived QSE & Use operators inherited from the Paper-I ground-state ansatz sequence.\\
\bottomrule
\end{tabular}
\end{center}
This is why the benchmark is partly about methods and partly about useful excitation operators.  The mixed system needs mixed fermion--boson operators, but different operator sources can still be tested as baselines.

\section{Why the overlap matrix matters}
The overlap matrix \(S\) is the Gram matrix of the selected QSE directions.  If two basis vectors are nearly the same direction, then \(S\) has a small eigenvalue.  The condition estimate
\begin{equation}
  \kappa(S)=\frac{\lambda_{\max}(S)}{\lambda_{\min}(S)}
\end{equation}
measures how fragile the generalized eigenproblem is.  A large \(\kappa(S)\) means that small numerical, sampling, or measurement errors can move roots substantially.

Regular mixed Fullmeta QSE can be spectrally accurate in ideal arithmetic because it includes many useful directions.  The same overcomplete alphabet can make \(S\) nearly singular.  Geometry-selected QSE tries to keep directions that explain spectral response without keeping redundant directions.

\section{Metric novelty as a residual}
Let the already-selected projected QSE directions be the columns of
\begin{equation}
  T=[t_1\ \cdots\ t_m],
\end{equation}
and let \(x\) be a new candidate projected direction.  The old span explains \(x\) by choosing coefficients \(\beta\):
\begin{equation}
  T\beta \approx x.
\end{equation}
Choose \(\beta\) by least squares:
\begin{equation}
  \beta^\star \in \arg\min_\beta \|x-T\beta\|^2.
\end{equation}
Expanding gives
\begin{equation}
  \|x-T\beta\|^2
  = x^\dagger x -2\operatorname{Re}(\beta^\dagger T^\dagger x)
  +\beta^\dagger T^\dagger T\beta.
\end{equation}
Define
\begin{equation}
  G=T^\dagger T,
  \qquad q=T^\dagger x,
  \qquad F=x^\dagger x.
\end{equation}
The normal equation is
\begin{equation}
  G\beta^\star=q.
\end{equation}
After optimizing out \(\beta\), the unexplained fraction is
\begin{equation}
  \mathcal N(x\mid T)
  =1-\frac{q^\dagger G^+q}{F}.
\end{equation}
This is the metric novelty: it is large when the candidate points outside the old selected span.

\begin{paperbox}
The optimized coefficients \(\beta^\star\) do not appear in the reported novelty because the old span has already been allowed to explain as much of the candidate as possible.
\end{paperbox}

\section{Spectral intensity and spectral L2}
For a probe observable \(O\), the QSE spectral function is a broadened weighted root list:
\begin{equation}
  I_{\QSE}(\omega)
  =\sum_\nu w_\nu^{\QSE}K_\eta(\omega-\omega_\nu^{\QSE}),
\end{equation}
where
\begin{equation}
  \omega_\nu^{\QSE}=E_\nu^{\QSE}-E_{\rm ref},
  \qquad
  w_\nu^{\QSE}=|\matel{\Psi_\nu^{\QSE}}{O}{\psi_0}|^2.
\end{equation}
The same-cutoff exact spectrum \(I_{\ED}(\omega)\) is built the same way using exact eigenvectors in the encoded Hilbert space.  The spectral diagnostic is
\begin{equation}
  \epsilon_{\rm spec}
  =\left(\sum_j |I_{\QSE}(\omega_j)-I_{\ED}(\omega_j)|^2\Delta\omega\right)^{1/2}
	ag{diagnostic}
\end{equation}
up to the implementation's grid normalization convention.

A sorted energy ladder asks whether root \(\nu\) matches root \(\nu\).  A spectrum asks whether the bright weighted peaks are reproduced.  These are different tests.  Dark exact roots can make the ordered ladder look poor while the spectrum overlay looks good.

\section{Resource proxy}
The matrix proxy \(S_{\rm mat}^{\rm proxy}\) estimates the amount of QSE matrix information required.  It grows with the number of basis pairs and Hamiltonian terms.  For a basis of size \(M\), the number of upper-triangular basis pairs is
\begin{equation}
  \frac{M(M+1)}{2}.
\end{equation}
The Hamiltonian matrix measurement burden scales like this pair count times the Hamiltonian term count.  The proxy is not a physical runtime claim; it is a consistent diagnostic for comparing basis policies.

\section{Current weak--weak pilot numbers}
For the current Hubbard--Holstein weak--weak diagnostic point, \((U/t,\lambda)=(0.25,0.25)\), \(n_{\rm ph}=2\), the useful pilot table is:
\begin{center}
\scriptsize
\begin{tabular}{lrrrr}
\toprule
Basis & rank & \(\kappa(S)\) & \(S_{\rm mat}\) & L2\\
\midrule
regular Fullmeta & 36 & \(8.08{\times}10^9\) & 289640 & \(2.31{\times}10^{-8}\)\\
geometry Fullmeta & 23 & \(2.13{\times}10^3\) & 9407 & \(4.63{\times}10^{-3}\)\\
Hamiltonian terms & 21 & \(6.28{\times}10^9\) & 19107 & \(2.42{\times}10^{-2}\)\\
ground ansatz & 6 & \(4.54{\times}10^2\) & 959 & \(1.05{\times}10^{-1}\)\\
\bottomrule
\end{tabular}
\end{center}
The regular Fullmeta row is almost exact on the broadened spectrum, but the overlap matrix is extremely ill-conditioned.  The geometry-selected row gives a compact, better-conditioned compromise.  The Hamiltonian-term row is not enough to fix conditioning.  The ground-ansatz-derived row is cheap and stable, but it misses more spectral weight.

\section{Promoted plots for the pilot}
The current promoted plot set is deliberately small.
\begin{enumerate}
\item \textbf{Spectrum overlay:} exact ED versus the QSE basis constructions.
\item \textbf{Energy-gap ladder:} sorted excitation gaps, labeled as a diagnostic rather than a one-to-one physical assignment.
\item \textbf{Tradeoff plot:} spectral L2 versus \(\kappa(S)\) and spectral L2 versus \(S_{\rm mat}^{\rm proxy}\).
\end{enumerate}
Transition-strength root-index scatter plots are not promoted at this stage because the visual grammar is not yet readable.

\section{Next code questions}
For the weak--weak pilot, the immediate code tasks are:
\begin{enumerate}
\item lock the promoted plot PDF format;
\item add additional probe observables only after the row definitions stay fixed;
\item decide whether Table~1 should report one observable or a small observable suite;
\item then repeat the same benchmark structure for the other Hubbard--Holstein regimes.
\end{enumerate}

\end{document}
